% GENERATED FILE. Edit canonical lessons, then run npm run generate:books.
% canonical-source-hash: fnv1a64:5d53965987eb3151
% canonical-lessons: ML-C415-samuham, ML-C415-potujanam, ML-C415-janasamkhya, ML-C415-samatvam, ML-C415-svatantryam

\chapter{Society}
\label{ch:ml-society}
\hlchaptermodality{415}

\begin{chapteropening}
\textbf{By the end of this chapter:} \emph{I can say society, the public, population, equality and freedom.}

Complete the last lesson of chapter 415: I can say society, the public, population, equality and freedom.
\end{chapteropening}

\section[\texorpdfstring{\ml{സമൂഹം} (samūhaṁ)}{സമൂഹം (samūhaṁ)}]{\ml{സമൂഹം} (samūhaṁ) — society}

\label{lesson:ML-C415-samuham}



\begin{quote}
Before the new one: say the Malayalam for a source, then the Malayalam for per cent.
\end{quote}

\subsection*{You'll want to know: \ml{സമൂഹം}}
\textbf{\ml{സമൂഹം}} — \emph{samūhaṁ} — \textquotedblleft{}society\textquotedblright{}.

\begin{quote}
\ml{സമൂഹത്തിൽ} \ml{എല്ലാവർക്കും} \ml{ഒരേ} \ml{അവകാശമാണ്}.
\end{quote}

\begin{quote}
\emph{samūhattil ellāvarkkuṁ orē avakāśamāṇŭ.}
\end{quote}

\begin{quote}
Everyone has the same rights in society.
\end{quote}

The adjective is \ml{സാമൂഹിക}, social: \ml{സാമൂഹിക} \ml{പ്രശ്നം}, a social problem.

\begin{cousinweb}
\textbf{\ml{സമൂഹം}} (\emph{samūhaṁ}) is borrowed from Sanskrit \emph{samūha}, \textquotedblleft{}a heap, a collection, an assembly\textquotedblright{} (CDIAL 13241). A society is people gathered together.
\end{cousinweb}

\begin{grammarlens}[title={What you've built}]
One more word for giving reasons and opinions.
\end{grammarlens}

\subsection*{Your turn}
\begin{itemize}
\raggedright
  \item \emph{Say it:} \emph{samūhaṁ}
  \item \emph{Say it:} \emph{samūhaṁ}, once more
  \item \emph{Say it:} \emph{samūhattil ellāvarkkuṁ orē avakāśamāṇŭ.} — Everyone has the same rights in society.
  \item \emph{Recall:} say the Malayalam for a source, then the Malayalam for per cent, then say \emph{samūhaṁ} again
\end{itemize}

\begin{tcolorbox}[breakable,colback=teal!4,colframe=teal!35!black,title={Before you move on}]
What does \ml{സമൂഹം} mean? (\textquotedblleft{}Society\textquotedblright{}.) Say it once more.
\end{tcolorbox}
\section[\texorpdfstring{\ml{പൊതുജനം} (potujanaṁ)}{പൊതുജനം (potujanaṁ)}]{\ml{പൊതുജനം} (potujanaṁ) — the general public}

\label{lesson:ML-C415-potujanam}



\begin{quote}
Before the new one: say the Malayalam for per cent, then the Malayalam for society.
\end{quote}

\subsection*{You'll want to know: \ml{പൊതുജനം}}
\textbf{\ml{പൊതുജനം}} — \emph{potujanaṁ} — \textquotedblleft{}the general public\textquotedblright{}.

\begin{quote}
\ml{പാർക്ക്} \ml{പൊതുജനത്തിന്} \ml{തുറന്നുകൊടുത്തു}.
\end{quote}

\begin{quote}
\emph{pārkkŭ potujanattinŭ tuṟannukoṭuttu.}
\end{quote}

\begin{quote}
The park was opened to the public.
\end{quote}

Often in the plural: \ml{പൊതുജനങ്ങൾ}, members of the public.

\begin{cousinweb}
\textbf{\ml{പൊതുജനം}} (\emph{potujanaṁ}) is a Malayalam compound: \ml{പൊതു}, \textquotedblleft{}common, general\textquotedblright{}, native Dravidian (DEDR 4507), and \ml{ജനം}, \textquotedblleft{}people\textquotedblright{}, from Sanskrit \emph{jana}, \textquotedblleft{}a person, a people\textquotedblright{} (CDIAL 5098).
\end{cousinweb}

\begin{grammarlens}[title={What you've built}]
One more word for giving reasons and opinions.
\end{grammarlens}

\subsection*{Your turn}
\begin{itemize}
\raggedright
  \item \emph{Say it:} \emph{potujanaṁ}
  \item \emph{Say it:} \emph{potujanaṁ}, once more
  \item \emph{Say it:} \emph{pārkkŭ potujanattinŭ tuṟannukoṭuttu.} — The park was opened to the public.
  \item \emph{Recall:} say the Malayalam for per cent, then the Malayalam for society, then say \emph{potujanaṁ} again
\end{itemize}

\begin{tcolorbox}[breakable,colback=teal!4,colframe=teal!35!black,title={Before you move on}]
What does \ml{പൊതുജനം} mean? (\textquotedblleft{}The general public\textquotedblright{}.) Say it once more.
\end{tcolorbox}
\section[\texorpdfstring{\ml{ജനസംഖ്യ} (janasaṁkhya)}{ജനസംഖ്യ (janasaṁkhya)}]{\ml{ജനസംഖ്യ} (janasaṁkhya) — population}

\label{lesson:ML-C415-janasamkhya}



\begin{quote}
Before the new one: say the Malayalam for society, then the Malayalam for the general public.
\end{quote}

\subsection*{You'll want to know: \ml{ജനസംഖ്യ}}
\textbf{\ml{ജനസംഖ്യ}} — \emph{janasaṁkhya} — \textquotedblleft{}population\textquotedblright{}.

\begin{quote}
\ml{നഗരത്തിലെ} \ml{ജനസംഖ്യ} \ml{വർഷം} \ml{തോറും} \ml{കൂടുന്നു}.
\end{quote}

\begin{quote}
\emph{nagarattile janasaṁkhya varṣaṁ tōṟuṁ kūṭunnu.}
\end{quote}

\begin{quote}
The city's population grows every year.
\end{quote}

Population is counted, so it goes up, \ml{കൂടുക}, or down, \ml{കുറയുക}.

\begin{cousinweb}
\textbf{\ml{ജനസംഖ്യ}} (\emph{janasaṁkhya}) is built from Sanskrit \emph{jana}, \textquotedblleft{}people\textquotedblright{} (CDIAL 5098), and \emph{saṁkhyā}, \textquotedblleft{}number, counting\textquotedblright{} (CDIAL 12841): the number of the people.
\end{cousinweb}

\begin{grammarlens}[title={What you've built}]
One more word for giving reasons and opinions.
\end{grammarlens}

\subsection*{Your turn}
\begin{itemize}
\raggedright
  \item \emph{Say it:} \emph{janasaṁkhya}
  \item \emph{Say it:} \emph{janasaṁkhya}, once more
  \item \emph{Say it:} \emph{nagarattile janasaṁkhya varṣaṁ tōṟuṁ kūṭunnu.} — The city's population grows every year.
  \item \emph{Recall:} say the Malayalam for society, then the Malayalam for the general public, then say \emph{janasaṁkhya} again
\end{itemize}

\begin{tcolorbox}[breakable,colback=teal!4,colframe=teal!35!black,title={Before you move on}]
What does \ml{ജനസംഖ്യ} mean? (\textquotedblleft{}Population\textquotedblright{}.) Say it once more.
\end{tcolorbox}
\section[\texorpdfstring{\ml{സമത്വം} (samatvaṁ)}{സമത്വം (samatvaṁ)}]{\ml{സമത്വം} (samatvaṁ) — equality}

\label{lesson:ML-C415-samatvam}



\begin{quote}
Before the new one: say the Malayalam for the general public, then the Malayalam for population.
\end{quote}

\subsection*{You'll want to know: \ml{സമത്വം}}
\textbf{\ml{സമത്വം}} — \emph{samatvaṁ} — \textquotedblleft{}equality\textquotedblright{}.

\begin{quote}
\ml{സ്ത്രീക്കും} \ml{പുരുഷനും} \ml{സമത്വം} \ml{വേണം}.
\end{quote}

\begin{quote}
\emph{strīkkuṁ puruṣanuṁ samatvaṁ vēṇaṁ.}
\end{quote}

\begin{quote}
Women and men need equality.
\end{quote}

The -\ml{ത്വം} ending makes an abstract noun, like English -ness.

\begin{cousinweb}
\textbf{\ml{സമത്വം}} (\emph{samatvaṁ}) is borrowed from Sanskrit \emph{samatva}, \textquotedblleft{}equality\textquotedblright{} (Monier-Williams), \emph{sama}, \textquotedblleft{}same, equal\textquotedblright{}, with the abstract ending \emph{-tva}.
\end{cousinweb}

\begin{grammarlens}[title={What you've built}]
One more word for giving reasons and opinions.
\end{grammarlens}

\subsection*{Your turn}
\begin{itemize}
\raggedright
  \item \emph{Say it:} \emph{samatvaṁ}
  \item \emph{Say it:} \emph{samatvaṁ}, once more
  \item \emph{Say it:} \emph{strīkkuṁ puruṣanuṁ samatvaṁ vēṇaṁ.} — Women and men need equality.
  \item \emph{Recall:} say the Malayalam for the general public, then the Malayalam for population, then say \emph{samatvaṁ} again
\end{itemize}

\begin{tcolorbox}[breakable,colback=teal!4,colframe=teal!35!black,title={Before you move on}]
What does \ml{സമത്വം} mean? (\textquotedblleft{}Equality\textquotedblright{}.) Say it once more.
\end{tcolorbox}
\section[\texorpdfstring{\ml{സ്വാതന്ത്ര്യം} (svātantryaṁ)}{സ്വാതന്ത്ര്യം (svātantryaṁ)}]{\ml{സ്വാതന്ത്ര്യം} (svātantryaṁ) — freedom}

\label{lesson:ML-C415-svatantryam}



\begin{quote}
Before the new one: say the Malayalam for population, then the Malayalam for equality.
\end{quote}

\subsection*{You'll want to know: \ml{സ്വാതന്ത്ര്യം}}
\textbf{\ml{സ്വാതന്ത്ര്യം}} — \emph{svātantryaṁ} — \textquotedblleft{}freedom\textquotedblright{}.

\begin{quote}
\ml{കുട്ടികൾക്ക്} \ml{കളിക്കാൻ} \ml{കുറച്ച്} \ml{സ്വാതന്ത്ര്യം} \ml{വേണം}.
\end{quote}

\begin{quote}
\emph{kuṭṭikaḷkkŭ kaḷikkān kuṟaccŭ svātantryaṁ vēṇaṁ.}
\end{quote}

\begin{quote}
Children need some freedom to play.
\end{quote}

A free country: \ml{സ്വതന്ത്ര} \ml{രാജ്യം}.

\begin{cousinweb}
\textbf{\ml{സ്വാതന്ത്ര്യം}} (\emph{svātantryaṁ}) is borrowed from Sanskrit \emph{svātantrya}, \textquotedblleft{}following one's own will, freedom\textquotedblright{} (Monier-Williams), made from \emph{svatantra}, \textquotedblleft{}under one's own rule\textquotedblright{}.
\end{cousinweb}

\begin{grammarlens}[title={What you've built}]
One more word for giving reasons and opinions.
\end{grammarlens}

\subsection*{Your turn}
\begin{itemize}
\raggedright
  \item \emph{Say it:} \emph{svātantryaṁ}
  \item \emph{Say it:} \emph{svātantryaṁ}, once more
  \item \emph{Say it:} \emph{kuṭṭikaḷkkŭ kaḷikkān kuṟaccŭ svātantryaṁ vēṇaṁ.} — Children need some freedom to play.
  \item \emph{Recall:} say the Malayalam for population, then the Malayalam for equality, then say \emph{svātantryaṁ} again
\end{itemize}

\begin{tcolorbox}[breakable,colback=teal!4,colframe=teal!35!black,title={Before you move on}]
What does \ml{സ്വാതന്ത്ര്യം} mean? (\textquotedblleft{}Freedom\textquotedblright{}.) Say it once more.
\end{tcolorbox}
